% Part: computability
% Chapter: recursive-functions
% Section: sequences

\documentclass[../../../include/open-logic-section]{subfiles}

\begin{document}

\olfileid[fr]{cmp}{rec}{seq}
\olsection{Suites finies}

L'ensemble des fonctions récursives primitives possède de
remarquables propriétés de clôture. Nous pourrons en obtenir
d'autres dès que nous disposerons d'un moyen convenable de
manipuler les \emph{suites finies}. Nous coderons une suite
finie d'entiers naturels par un entier naturel de la manière
suivante : à
$\langle a_0,a_1,a_2,\dots,a_k\rangle$, nous associons
\[
p_0^{a_0+1} \cdot p_1^{a_1+1} \cdot p_2^{a_2+1} \cdot \dots \cdot
p_k^{a_k+1}.
\]
Nous ajoutons $1$ à chaque exposant afin que, par exemple,
les suites $\langle 2,7,3\rangle$ et
$\langle 2,7,3,0,0\rangle$ aient des codes distincts.
Les nombres $0$ et~$1$ peuvent tous deux servir de codes à
la suite vide; pour fixer la notation, $\emptyseq$
désignera~$0$.

Ce codage repose sur le théorème fondamental de
l'arithmétique : tout entier naturel $n\geq2$ possède une
unique écriture de la forme
\[
n = p_0^{a_0} \cdot p_1^{a_1} \cdot \dots \cdot p_k^{a_k}
\]
où les $a_i$ sont des entiers naturels et $a_k\geq1$.
L'application
$\tuple{}(a_0,\dots,a_k)=\tuple{a_0,\dots,a_k}$
est donc injective : des suites différentes reçoivent des
codes différents, et chaque entier représente au plus une
suite dans ce codage.

Montrons maintenant que l'on peut calculer par des fonctions
récursives primitives la longueur d'une suite, son élément
de rang~$i$, l'adjonction d'un élément à sa fin et la
concaténation de deux suites.

\begin{prop}
  La fonction $\len{s}$, qui renvoie la longueur de la suite
  codée par~$s$, est récursive primitive.
\end{prop}

\begin{proof}
  Soit $R(i,s)$ la relation définie par
  \[
  R(i, s) \text{ si et seulement si }
  p_i \mid s \land p_{i+1} \nmid s.
  \]
  Elle est récursive primitive. Si $s$ code une suite non vide,
  c'est-à-dire si
  \[
  s = p_0^{a_0+1} \cdot \dots \cdot p_{k}^{a_{k}+1},
  \]
  alors $R(i,s)$ vaut exactement lorsque $p_i$ est le plus
  grand nombre premier divisant~$s$, autrement dit lorsque
  $i=k$. La longueur vaut donc $i+1$ dans ce cas. Posons
  \[
  \len{s} =
  \begin{cases}
    0 & \text{si $s = 0$ ou $s = 1$} \\
    1 + \bmin{i < s}{R(i, s)} & \text{sinon}
  \end{cases}
  \]
  La minimisation est bornée : pour tout code non vide,
  il existe un unique~$i$ vérifiant $R(i,s)$, et cet
  indice est strictement inférieur à~$s$.
\end{proof}

\begin{prop}
  La fonction $\fn{append}(s,a)$, qui renvoie le code obtenu
  en ajoutant~$a$ à la fin de la suite codée par~$s$,
  est récursive primitive.
\end{prop}

\begin{proof}
  Elle se définit par
  \[
  \fn{append}(s,a) =
  \begin{cases}
    2^{a+1} & \text{si $s = 0$ ou $s = 1$} \\
    s \cdot p_{\len{s}}^{a+1} & \text{sinon.}
  \end{cases}
  \]
\end{proof}

\begin{prop}
  La fonction $\fn{element}(s,i)$, qui renvoie l'élément
  de rang~$i$ de la suite codée par~$s$ (le premier
  élément est de rang~$0$), ou~$0$ si
  $i\geq\len{s}$, est récursive primitive.
\end{prop}

\begin{proof}
  L'entier~$a$ est l'élément de rang~$i$ de~$s$
  si et seulement si $p_i^{a+1}$ est la plus grande
  puissance de~$p_i$ qui divise~$s$ : on a
  $p_i^{a+1}\mid s$ et $p_i^{a+2}\nmid s$. Ainsi,
  \[
  \fn{element}(s,i) =
  \begin{cases}
    0 & \mbox{si $i \geq \len{s}$} \\
    \bmin{a < s}{(p_i^{a+2} \nmid s)} & \text{sinon.}
  \end{cases}
  \]
\end{proof}

Nous emploierons désormais une notation plus courte :
$(s)_i$ pour $\fn{element}(s,i)$, et
$\tuple{s_0,\dots,s_k}$ pour
\[
\fn{append}(\fn{append}(\dots \fn{append}(\emptyseq,s_0)
\dots),s_k).
\]
Si la suite~$s$ a pour longueur~$k$, ses éléments sont
$(s)_0,\dots,(s)_{k-1}$.

\begin{prop}
  La fonction $\fn{concat}(s,t)$, qui concatène les
  deux suites codées par $s$ et~$t$, est récursive primitive.
\end{prop}

\begin{proof}
  Nous cherchons une fonction $\fn{concat}$ telle que
  \[
    \fn{concat}(\tuple{a_0, \dots, a_k}, \tuple{b_0, \dots, b_l}) =
    \tuple{a_0, \dots, a_k, b_0, \dots, b_l}.
  \]
  Définissons d'abord une fonction auxiliaire
  $\fn{hconcat}(s,t,n)$ qui adjoint à~$s$ les $n$
  premiers éléments de~$t$. La récurrence primitive donne
  \begin{align*}
    \fn{hconcat}(s,t,0) & = s\\
    \fn{hconcat}(s,t,n+1) & = \fn{append}(\fn{hconcat}(s,t,n),(t)_n)
    \intertext{Nous pouvons alors définir $\fn{concat}$ par}
    \fn{concat}(s,t) & = \fn{hconcat}(s,t,\len{t}).
  \end{align*}
\end{proof}

Nous écrirons $s\concat t$ pour $\fn{concat}(s,t)$.

Pour utiliser la recherche bornée, il nous faut une borne
sur le code d'une suite en fonction de sa longueur et de
son plus grand élément. Supposons que $s$ code une suite
de longueur~$k$ dont chaque élément est inférieur ou égal
à~$x$. Son code a au plus $k$ facteurs premiers, chacun
inférieur ou égal à~$p_{k-1}$ et élevé à une puissance
au plus égale à~$x+1$. Pour $k\geq1$, posons
\[
\fn{sequenceBound}(x,k) = p_{k-1}^{k \cdot (x+1)},
\]
et, pour la suite vide, posons
$\fn{sequenceBound}(x,0)=1$. Le code de la suite~$s$
est alors inférieur ou égal à cette borne. La valeur
pour $k=0$ évite l'indice $p_{-1}$ et se calcule par
une définition par cas récursive primitive.\footnote{Complément éditorial : la source ne définit pas la borne lorsque $k=0$, où $p_{k-1}$ n'a pas de sens. La valeur $1$ couvre les deux codes admis de la suite vide.}

Cette borne permet de définir de nouvelles fonctions
par recherche bornée. Par exemple, on peut redéfinir
$\fn{concat}$ : il suffit de donner une
\emph{spécification} récursive primitive du code
cherché et une borne de recherche. Pour les codes
canoniques produits à partir de $\emptyseq=0$, on obtient le schéma ci-dessous. Il faut aussi exiger que le candidat $v$ soit un code de suite. Écrivons $\fn{SeqCode}(v)$ pour la relation $v\leq1\lor(v\geq2\land\bforall{i<v}{(p_{i+1}\mid v\lif p_i\mid v)})$. Cette relation récursive primitive exige une chaîne de facteurs premiers sans lacune :
\begin{align*}
  \fn{concat}(s,t) = {} & \bmin{v < \fn{sequenceBound}(s+t,\len{s} +
    \len{t})}{} \\
  & \quad(\fn{SeqCode}(v) \land \len{v} = \len{s} + \len{t} \land {}\\
    & \qquad \bforall{i < \len{s}}{(v)_i = (s)_i} \land {} \\
    & \qquad \bforall{j < \len{t}}{(v)_{\len{s}+j} = (t)_j})
\end{align*}
Le code voulu est strictement sous cette borne : pour une
suite non vide, chaque élément est strictement inférieur à $s+t$, donc chaque exposant est au plus $s+t$ ; le produit est strictement sous la borne.
Pour deux suites vides, le code choisi est $0<1$.\footnote{Corrections éditoriales : les deux quantifications portant sur les segments forment deux conditions distinctes dans une conjonction ; dans la source, la seconde est placée sous la première et disparaît quand le premier segment est vide. Nous ajoutons aussi le test de code : sans lui, pour $s=t=2$, la recherche trouve $v=3$ au lieu du code $6$ de la suite $\tuple{0,0}$.}
Si l'on emploie aussi~$1$ comme code de la suite vide,
les deux procédés peuvent rendre respectivement $1$ et~$0$
pour deux suites vides : ces codes représentent la même suite,
mais les fonctions numériques ne coïncident alors pas partout.

\begin{prob}
Montrer qu'il existe une fonction récursive primitive
$\fn{sconcat}(s)$ telle que
\[
\fn{sconcat}(\tuple{s_0, \dots, s_k}) = s_0 \concat \dots \concat s_k.
\]
\end{prob}

\begin{prob}
Montrer qu'il existe une fonction récursive primitive
$\fn{tail}(s)$ telle que
\begin{align*}
  \fn{tail}(\emptyseq) & = 0 \text{ et}\\
  \fn{tail}(\tuple{s_0, \dots, s_{k}}) & = \tuple{s_1, \dots, s_{k}}.
\end{align*}
Pour $k=0$, le membre de droite désigne la suite vide.
\end{prob}

\begin{prop}
  \ollabel{prop:subseq}
  La fonction $\fn{subseq}(s,i,n)$, qui renvoie la
  sous-suite de~$s$ de longueur~$n$ commençant à
  l'élément de rang~$i$, est récursive primitive.
  Cette description porte sur $i+n\leq\len{s}$;
  hors de ce domaine, on peut convenir qu'elle renvoie
  le code de la suite vide.\footnote{Complément éditorial : la source ne précise pas le cas où une sous-suite de cette longueur n'existe pas. Nous fixons ce retour pour donner une fonction totale.}
\end{prop}

\begin{proof}
  Exercice.
\end{proof}

\begin{prob}
  Démontrer la \olref[cmp][rec][seq]{prop:subseq}.
\end{prob}

\end{document}
